\documentclass[11pt,a4paper]{article}
\usepackage[spanish]{babel}
\usepackage[utf8]{inputenc}
\usepackage[T1]{fontenc}
\usepackage{lmodern}
\usepackage[margin=2.35cm]{geometry}
\usepackage{graphicx,xcolor,listings,float,fancyhdr,booktabs,hyperref}
\definecolor{azul}{HTML}{173A63}
\definecolor{suave}{HTML}{F1F5F9}
\hypersetup{colorlinks=true,urlcolor=azul,linkcolor=azul}
\lstset{language=Python,basicstyle=\ttfamily\small,breaklines=true,frame=single,columns=fullflexible,showstringspaces=false,keywordstyle=\color{azul},backgroundcolor=\color{suave}}
\pagestyle{fancy}\fancyhf{}\lhead{SRI -- Mayusculator UDP}\rhead{Python 3}\cfoot{\thepage}\setlength{\headheight}{15pt}
\setlength{\parindent}{0pt}\setlength{\parskip}{5pt}
% Sustituir los datos antes de entregar.
\newcommand{\AlumnoUno}{Nombre y apellidos del alumno 1}
\newcommand{\AlumnoDos}{Nombre y apellidos del alumno 2}
\newcommand{\Repositorio}{https://github.com/KarmaVS/Mayusculator-UDP-Python3}
\newcommand{\captura}[2]{%
\begin{figure}[H]\centering
\IfFileExists{codigo/imagenes/#1}{\includegraphics[width=.96\linewidth,height=5.4cm,keepaspectratio]{codigo/imagenes/#1}}{%
\fcolorbox{azul}{suave}{\parbox[c][4.7cm][c]{.91\linewidth}{\centering\textcolor{azul}{\bfseries CAPTURA PENDIENTE}\\[6pt]\texttt{\detokenize{codigo/imagenes/#1}}\\[5pt]Añadir la imagen real antes de entregar.}}}
\caption{#2}\end{figure}}
\begin{document}
\begin{titlepage}
\centering
\vspace*{2.7cm}
{\Large Sistemas de Redes e Internet (SRI)\par}
\vspace{1.4cm}
{\color{azul}\Huge\bfseries Mayusculator\par}
\vspace{.45cm}
{\LARGE Comunicación cliente-servidor UDP\par}
\vspace{.5cm}{\large Adaptación a Python 3 y pruebas de red\par}
\par\vspace{1cm}\rule{.75\linewidth}{.5pt}\par\vspace{1cm}
{\large \AlumnoUno\\[.3cm]\AlumnoDos\par}
\vfill
{\large Servidor: Debian\\Cliente: Ubuntu\\IP del servidor: 192.168.3.3\\Puerto: 12000/UDP\par}
\vspace{1cm}{\large Curso 2026--2027\par}
\end{titlepage}
\tableofcontents
\vfill
\textbf{Nota de preparación.} Esta plantilla documenta la versión UDP. Las capturas y los resultados están pendientes de las pruebas reales; sustituir los datos de los autores y confirmar el enlace del repositorio antes de entregar.
\newpage
\section{Introducción y objetivos}
Mayusculator es una aplicación cliente-servidor que permite enviar una cadena de texto a través de la red. El servidor convierte la cadena a mayúsculas mediante \texttt{upper()} y devuelve el resultado al cliente. En esta práctica se utiliza UDP y Python 3.
\begin{itemize}
\item Comprender la comunicación mediante datagramas UDP.
\item Adaptar los programas para enviar y recibir bytes en Python 3.
\item Comprobar la comunicación entre dos máquinas distintas.
\item Documentar la configuración y las pruebas con capturas reales.
\end{itemize}
\section{Funcionamiento de UDP}
UDP (User Datagram Protocol) es un protocolo de transporte no orientado a conexión. No establece una sesión previa ni garantiza por sí mismo la entrega, el orden o la ausencia de duplicados. En esta práctica el cliente envía un datagrama al servidor y espera otro datagrama con la respuesta.
\begin{center}
\begin{tabular}{lll}\toprule
Característica & TCP & UDP\\\midrule
Tipo de socket & \texttt{SOCK\_STREAM} & \texttt{SOCK\_DGRAM}\\
Conexión previa & Sí & No\\
Envío & \texttt{sendall()} & \texttt{sendto()}\\
Recepción & \texttt{recv()} & \texttt{recvfrom()}\\
Entrega y orden fiables & Sí, si no falla la conexión & No\\\bottomrule
\end{tabular}
\end{center}
\section{Entorno de trabajo}
Se propone reutilizar el entorno de la práctica TCP: servidor Debian, cliente Ubuntu y red IPv4. Se debe confirmar que la configuración sigue siendo válida antes de ejecutar las pruebas.
\begin{center}\begin{tabular}{ll}\toprule
Parámetro & Valor previsto\\\midrule
IP del servidor & 192.168.3.3/24\\
Interfaz del servidor & enp0s3\\
Puerto & 12000/UDP\\
Python del servidor & 3.13.5 (comprobar)\\
Cliente & Ubuntu, IP propia de su equipo\\\bottomrule\end{tabular}\end{center}
\newpage
\section{Código y adaptación a Python 3}
\subsection{Servidor UDP}
Se crea un socket IPv4 de tipo datagrama y se asocia al puerto 12000. La función \texttt{recvfrom()} devuelve los bytes recibidos y la dirección de origen. Tras convertir el texto a mayúsculas se responde mediante \texttt{sendto()}.
\lstinputlisting{codigo/servidorudp.py}
\subsection{Cliente UDP}
El cliente solicita una cadena al usuario, envía un datagrama a la IP del servidor y espera el datagrama de respuesta. Al terminar, cierra el socket.
\lstinputlisting{codigo/clienteudp.py}
\subsection{Cambios principales}
\texttt{input()} permite introducir texto. \texttt{encode("utf-8")} convierte cadenas en bytes y \texttt{decode("utf-8")} realiza la operación inversa. El límite de recepción es de 1024 bytes por datagrama, adecuado para mensajes cortos. No se necesita \texttt{listen()}, \texttt{accept()} ni \texttt{connect()} en esta implementación.
\newpage
\section{Configuración y comprobación de la red}
\subsection{Dirección IP en Debian}
En la máquina servidora se comprueba la IP mediante \texttt{ip a} y la versión de Python mediante \texttt{python3 --version}. Debe aparecer la dirección prevista 192.168.3.3/24. Si ya está configurada, no es necesario modificarla.
\captura{01_ip_debian.png}{Dirección IP y configuración del servidor Debian.}
\subsection{Conectividad desde Ubuntu}
Desde el cliente se ejecuta el siguiente comando para verificar la conectividad IP. Un ping correcto no demuestra por sí solo que el servicio UDP funcione.
\begin{lstlisting}[language=bash]
ping -c 4 192.168.3.3
\end{lstlisting}
\captura{02_conectividad.png}{Prueba de conectividad desde Ubuntu.}
\newpage
\section{Ejecución y pruebas UDP}
\subsection{Inicio del servidor}
En Debian, desde el directorio que contiene el código, se ejecuta:
\begin{lstlisting}[language=bash]
python3 servidorudp.py
\end{lstlisting}
El programa debe quedar esperando datagramas.
\captura{03_servidor_esperando.png}{Servidor UDP iniciado en Debian.}
\subsection{Comprobación del puerto}
En otra terminal de Debian se verifica que el socket UDP está asociado al puerto 12000. El estado habitual es UNCONN, no LISTEN.
\begin{lstlisting}[language=bash]
sudo ss -ulnp | grep ':12000'
\end{lstlisting}
\captura{04_puerto_udp.png}{Comprobación del socket UDP en el puerto 12000.}
\newpage
\subsection{Envío del mensaje desde Ubuntu}
En el equipo cliente, desde el directorio del programa, se ejecuta:
\begin{lstlisting}[language=bash]
python3 clienteudp.py
\end{lstlisting}
Como prueba se propone introducir \texttt{hola mundo}. La respuesta esperada es \texttt{HOLA MUNDO}.
\captura{05_cliente_ubuntu.png}{Envío del mensaje y respuesta recibida en Ubuntu.}
\subsection{Recepción en Debian}
Se vuelve a la terminal del servidor y se comprueba que aparece el mensaje original recibido. El código lo transforma y envía el resultado a la dirección del cliente.
\captura{06_servidor_recibiendo.png}{Mensaje recibido por el servidor Debian.}
\newpage
\section{Resultados y resolución de problemas}
Completar la tabla después de realizar las pruebas. Los resultados esperados no constituyen evidencia de que las pruebas se hayan realizado.
\begin{center}\begin{tabular}{lll}\toprule
Prueba & Resultado esperado & Resultado real\\\midrule
Dirección IP & 192.168.3.3/24 & Pendiente\\
Puerto UDP & 12000 asociado & Pendiente\\
Mensaje de prueba & HOLA MUNDO & Pendiente\\
Terminal del servidor & Mensaje recibido & Pendiente\\\bottomrule\end{tabular}\end{center}
Si el cliente queda esperando, se comprueba que el servidor esté ejecutándose, que la IP y el puerto coincidan, y que el cortafuegos permita UDP/12000 desde el cliente. Si UFW está activo, se puede autorizar únicamente la IP del cliente mediante una regla específica. El programa sencillo no incorpora tiempo de espera, por lo que puede interrumpirse con Ctrl+C si no llega respuesta.
\section{Repositorio y referencias}
El código modificado y las capturas pueden consultarse en el repositorio de GitHub (confirmar la dirección antes de entregar):
\begin{center}\href{\Repositorio}{\textbf{Repositorio Mayusculator UDP - Python 3}}\end{center}
% Si se elige otro repositorio, cambiar la variable Repositorio al principio.
\begin{itemize}
\item Python: \url{https://docs.python.org/es/3/library/socket.html}
\item Guía de sockets: \url{https://docs.python.org/3/howto/sockets.html}
\item RFC 768 (UDP): \url{https://www.rfc-editor.org/rfc/rfc768}
\end{itemize}
\section{Conclusión}
La práctica permite observar cómo dos programas pueden intercambiar datagramas sin establecer una conexión previa. A diferencia de TCP, UDP no confirma la entrega; por ello, una respuesta correcta del cliente y la recepción registrada por el servidor son las evidencias principales de funcionamiento. Completar esta conclusión con los resultados reales de la prueba.
\end{document}
